%==============================================================================
%  04-application-categories.tex — Section 3: AI Application Categories
%==============================================================================
\strategyPart{08 \textbar\ AI Technologies}{AI Applications That Solve Organizational Problems}

The correct unit of analysis is not the algorithm; it is the organizational
problem. Drawing on the problem-centric categorization used in MIT Sloan
research and McKinsey's use-case economics, this section groups AI
applications into four categories — each defined by the class of business
problem it solves [MIT SMR 2019; McKinsey 2023].

\section{Four problem-centric categories}

\vspace{-0.3em}
\aiCard{%
  \aiCardTitle{Predictive AI}\\[0.4em]
  \textit{Forecasting \& risk.} Demand and revenue forecasting, churn
  prediction, credit risk, predictive maintenance.
}\par\vspace{0.7em}

\aiCard{%
  \aiCardTitle{Optimization AI}\\[0.4em]
  \textit{Resource allocation.} Supply-chain planning, pricing and revenue
  management, scheduling, routing.
}\par\vspace{0.7em}

\aiCard{%
  \aiCardTitle{Generative \& Conversational AI}\\[0.4em]
  \textit{Knowledge \& creativity.} Document synthesis, semantic search,
  marketing content, copilots and agents.
}\par\vspace{0.7em}

\aiCard{%
  \aiCardTitle{Computer Vision}\\[0.4em]
  \textit{Physical-world interpretation.} Visual quality inspection, safety
  monitoring, document capture, autonomous operations.
}\par

\section{Where the value concentrates}

McKinsey's analysis of generative AI found that roughly \textbf{75\% of its
potential annual value} concentrates in four functions: customer operations,
marketing and sales, software engineering, and research and development
[McKinsey 2023]. Predictive and optimization AI remain the most widely
deployed capabilities in production, because they map directly to existing
P\&L line items — revenue, cost, risk and cash. The executive implication is
that a problem-first portfolio is not an even spread across all four
categories; it is a concentration on the categories that map to the
organization's most expensive unsolved problems.

\section{The AI technology stack}

\begin{center}
  \renewcommand{\arraystretch}{1.45}
  \rowcolors{2}{inSnow}{white}
  \fitwidth{\begin{tabular}{>{\raggedright\arraybackslash}p{2.6cm}>{\raggedright\arraybackslash}p{3.4cm}>{\raggedright\arraybackslash}p{2.4cm}>{\raggedright\arraybackslash}p{2.4cm}>{\raggedright\arraybackslash}p{2.4cm}}
    \rowcolor{inNavy}
    \hcell{Technology} & \hcell{Primary use} & \hcell{Autonomy} & \hcell{Human role} & \hcell{Main risk} \\
    \textbf{Predictive AI} & Forecasting, risk & Low & Approves & Forecast error \\
    \textbf{Classification / recommend} & Segmentation, next-best-action & Low--medium & Reviews & Bias \\
    \textbf{Generative AI / LLM} & Content, knowledge & Low & Verifies & Hallucination \\
    \textbf{RAG} & Grounded knowledge & Medium & Validates & Stale data \\
    \textbf{Copilots / assistants} & Augmented work & Low--medium & Reviewer & Overreliance \\
    \textbf{Autonomous agents} & Task execution & High & Supervisor & Unauthorized action \\
    \textbf{Multi-agent BI} & Coordinated analytics & High & Governor & Cascading failure \\
  \end{tabular}}
\end{center}

\section{Selecting the right category}

\begin{itemize}
  \item \textbf{Uncertainty about the future} $\rightarrow$ \textbf{Predictive AI}.
  \item \textbf{Allocation of scarce resources} $\rightarrow$ \textbf{Optimization AI}.
  \item \textbf{Known answers buried in unstructured knowledge} $\rightarrow$ \textbf{Generative \& Conversational AI}.
  \item \textbf{Signal trapped in images, video or physical assets} $\rightarrow$ \textbf{Computer Vision}.
\end{itemize}

Selection discipline is the first test of the problem-first philosophy: a
technology-first organization buys the shiniest model; a problem-first
organization buys the capability that retires the most expensive problem.

\section{Seven architecture patterns}

Anthropic's agent research recommends \textbf{selecting the simplest
architecture that reliably solves the business problem} — fixed workflows for
predictable, clearly-defined processes; agents only where dynamic interaction
is genuinely required [Anthropic]. The seven patterns:

\begin{enumerate}
  \item \textbf{Simple AI assistant} — employee $\rightarrow$ model
        $\rightarrow$ answer; suited to summarization, writing, translation
        and basic search.
  \item \textbf{Retrieval-augmented assistant} — search, retrieve, generate
        with citations; suited to policy, documentation, support and
        contract analysis.
  \item \textbf{Copilot workflow} — recommendation $\rightarrow$ human
        approval $\rightarrow$ action; suited to sales, service, product,
        finance and HR decisions.
  \item \textbf{Fixed AI workflow} — input $\rightarrow$ steps $\rightarrow$
        validation $\rightarrow$ output; suited to invoice processing,
        classification, standard reporting and ticket routing
        [Anthropic].
  \item \textbf{Single AI agent} — goal $\rightarrow$ planning $\rightarrow$
        tool use $\rightarrow$ result $\rightarrow$ evaluation; suited to
        research, IT troubleshooting and complex document analysis.
  \item \textbf{Multi-agent system} — orchestrator + specialists +
        validation + human approval; suited to complex BI, supply-chain,
        R\&D and cross-functional analysis.
  \item \textbf{Agent-to-agent organization} — manager agent coordinates
        finance, sales and operations agents under human executive decision;
        requires strong identity, access control and validation.
\end{enumerate}

The design principle: \textbf{assistants for judgment-heavy work, workflows
for predictable processes, agents for bounded and measurable tasks, and
multi-agent systems only when the problem genuinely requires dynamic
coordination} [Anthropic; Microsoft agent framework; NIST AI RMF].

\takeaway{Choose the AI category by the business problem (predictive, optimization, generative, vision) and use the simplest architecture pattern.}
